\documentclass[a4paper]{article}
\usepackage[pdftex]{graphicx}
\usepackage[utf8]{inputenc}
\usepackage{enumerate}
\usepackage{eurosym}
\usepackage{amssymb}
\usepackage{icomma}
\usepackage{siunitx}
\sisetup{locale=DE} 
\usepackage{tikz}
\usepackage{href-ul}
\hypersetup{
	colorlinks=true,
	linkcolor=blue,
	urlcolor=blue}
\usepackage{geometry}
\geometry{a4paper, top=15mm, left=15mm, right=15mm, bottom=15mm,
	headsep=10mm, footskip=12mm}

\begin{document}
\begin{tikzpicture}(10,3)
	\draw[ultra thick](2,0) --(10,0) -- (10,1.3) --(2,1.3) -- (2,0);
	\draw[fill=black](2,0)-- (0,.6) -- (2,.6) -- (2,0);
	\node at (6,.6) {\large Lösungsblatt von \href{https://www.okuyakl.de}{www.okuyakl.de}};
\end{tikzpicture}
\vspace{0.5 cm}

\noindent{\bf Aufgabe 1. a)}\\
Die Anzahl der Handbälle sei $x$. Dann ist die Anzahl der Fußbälle $2x$ und die der Volleybälle
$x+4$. Die 12 Tennisbälle kosten 12 $\cdot$ \EUR{6} = \EUR{72}. Die Gesamtkosten sind:
$$
\begin{array}{rcl}
K(x)&=&(x\cdot 30+2x \cdot 35 +(x+4)\cdot 20 +72)\,\textnormal{\EUR{}}\\
&=& (30x +70x +20x + 80 + 72)\, \textnormal{\EUR{}}\\
&=& (120x + 152)\, \textnormal{\EUR{}} 
\end{array}
$$
\vspace{0.5 cm}

\noindent{\bf Aufgabe 1. b)}\\
Nun setzen wir die Variable x = 10:
$$K(10)=120\cdot 10 + 152 = \textnormal{\EUR{1352}}$$
Die Gesamtkosten für die Bälle belaufen sich auf \EUR{1352}.
\vspace{0.5 cm}

\noindent{\bf Aufgabe 2.}\\
	Jede Breite (=x), Höhe (=2x), Länge (=4x) wird viermal abgewickelt. Die Schnurlänge $l$ ist:
	$$l=4x + 4\cdot 2x + 4 \cdot 4x = 28x$$
\vspace{0.5 cm}

\noindent{\bf Aufgabe 3. a) }\\
$$
\renewcommand{\arraystretch}{2}
\begin{array}{rclll}

\textnormal{1. Tag}&=&x~\SI{}{\kilo\meter}  \\
\textnormal{2. Tag}&=&2x~\SI{}{\kilo\meter}  \\
\textnormal{1. Tag}&=&2x-15~\SI{}{\kilo\meter}  \\
\hline
\textnormal{gesamt}&=&x+2x+2x-15&=5x-15\\
\textnormal{b)}\quad 235 &=&x+2x+2x-15&=5x-15 &|+15\\
250 &=& 5x &|:5\\
50 &=& x
\end{array}
$$

Er ist am ersten Tag $\SI{50}{\kilo\meter}$ gefahren.
\vspace{0.5 cm}

\noindent{\bf Aufgabe 4. a)}\\
$$
\renewcommand{\arraystretch}{2}
\begin{array}{rcll}
\textnormal{1. Mai}&=&\SI{100}{\gram}  \\
\textnormal{15. Mai}&=& \SI{200}{\gram} \\
\textnormal{29. Mai}&=& \SI{400}{\gram} \\
\hline
14&\cdot &\SI{0,1}{\kilo\gram} \\
14&\cdot &\SI{0,2}{\kilo\gram} \\
3&\cdot &\SI{0,4}{\kilo\gram} \\
\hline
&=& 1,4+2,8+1,2&=\SI{5,4}{\kilo\gram} \\
\end{array}
$$

Am 31. Mai bekommt der Hund $\SI{0,4}{\kilo\gram}$ Futter zu fressen.
b) Sie hat bis dahin 6 Futtertüten angebrochen.
\newpage

\noindent{\bf Aufgabe 5. a)}\\
$$
\renewcommand{\arraystretch}{2}
\begin{array}{rcll}
\textnormal{2 Erwachsene}&=&2x \\
\textnormal{3 Kinder}&=&3\cdot 0,5x\\
\textnormal{1 Auto}&=&9,50\\
\hline
20,70 &=& 3,5x+9,50 &|-9,50\\
11,20 &=&3,5x &|:3,5\\
3,20 &=&x
\end{array}
$$

Frau Jagmichel (1 Erwachsene) müsste \EUR{3,20} zahlen.
\vspace{0.5 cm}

\noindent{\bf Aufgabe 5. b)}\\
$$2 \cdot 3,20 \cdot 0,75 = 4,80$$
Die Großeltern würden noch  \EUR{4,80} kosten.
\vspace{0.5 cm}

\noindent{\bf Aufgabe 6.}\\
\begin{center}
	\renewcommand{\arraystretch}{2}
	\begin{tabular}{|c|p{1.5 cm}|p{1.5 cm}|p{1.5 cm}|p{1.5 cm}|}
		\hline
		Geschwindigkeit in km/h & 20 & 50 & 80& 100\\
		\hline
		Bremsweg in  m	& 4 & 25 & 64 & 100 \\
		\hline
	\end{tabular}
\end{center}
\vspace{0.5 cm}

\noindent{\bf Aufgabe 7. a)}\\
$$
\renewcommand{\arraystretch}{2}
\begin{array}{rcl}

\textnormal{1. Jahr}&=&12\cdot x \\
\textnormal{2. Jahr}&=&12\cdot (x+150)\\
\textnormal{3. Jahr}&=&12\cdot (x+300) +540\\
\hline
T(x)&=&36\cdot x + 12\cdot 150 +12\cdot 300 +540 \\
&=& 36x+5940\\
\end{array}
$$
\vspace{0.5 cm}


\begin{minipage}{0.5\textwidth}
	\noindent{\bf Aufgabe 7. b)}\\
	$$T(450)=36\cdot 450 + 5940 = 22140 $$
	Sein Gesamtverdienst wird \EUR{22140} betragen.
	\vspace{0.5 cm}
	
	\noindent{\bf Aufgabe 7. c)}\\
	$$22410 : 36 = 615$$
	Sein Durchschnittsverdienst wird \EUR{615} sein.
	\vspace{0.5 cm}
	
	\noindent{\bf Aufgabe 7. d)}\\
	$$22410+3\cdot 120 = 22500$$
	$$22500 : 36 = 625$$
	Sein Gesamtverdienst wird dann \EUR{22500} betragen.\\
	Sein Durchschnittsverdienst steigt auf \EUR{625}.
\end{minipage}
\hfill
\begin{minipage}{0.5\textwidth}
\begin{center}
	\includegraphics[width=7 cm]{../../viecher/endcomic.pdf}
	
	Hier geht es zurück zum \href{https://www.okuyakl.de/math/m7teaL118/aa118.pdf}{Aufgabenblatt}
\end{center}

\end{minipage}

\end{document}

